% ECONOMICS_PROBLEM: {"id": "I15-520", "title": "Mental-health poverty feedbacks", "jel_code": "I15", "source_id": "OP-0448"}
% !TeX program = xelatex
\documentclass[11pt,a4paper]{article}
\usepackage[margin=23mm]{geometry}
\usepackage{fontspec}
\setmainfont{Latin Modern Roman}
\usepackage{amsmath,amssymb,mathtools}
\usepackage{xcolor,enumitem,booktabs,longtable,array,xurl,hyperref}
\definecolor{muted}{HTML}{53636C}
\hypersetup{colorlinks=true,urlcolor=blue,linkcolor=blue}
\setlength{\parindent}{0pt}
\setlength{\parskip}{5pt}
\newcommand{\R}{\mathbb R}
\newcommand{\E}{\mathbb E}
\newcommand{\N}{\mathbb N}
\newcommand{\ind}{\mathbf 1}
\newcommand{\Var}{\operatorname{Var}}
\newcommand{\Cov}{\operatorname{Cov}}
\newcommand{\supp}{\operatorname{supp}}
\DeclareMathOperator*{\argmin}{arg\,min}
\DeclareMathOperator*{\argmax}{arg\,max}
\newcommand{\entrymeta}[3]{{\sffamily\small\color{muted}\textbf{\texttt{#1}}\quad|\quad #2\par #3\par}}
\newcommand{\Problemref}[1]{\texttt{#1}}
\newcommand{\JELref}[2]{\texttt{#2} (JEL \texttt{#1})}
\begin{document}
\section*{Mental-health poverty feedbacks}
\textbf{Problem ID:} \texttt{I15-520}\quad \textbf{JEL:} \texttt{I15}\par
% BEGIN ECONOMICS STATEMENT
\label{problem:OP-0448}
{\sffamily\small\color{muted}Original subject group: Health economics\par}
\label{definitions:problem:30007730}
\textbf{Audit classification:} Published research agenda. Checked source identity and appropriate agenda/background support; expanded intervention, units, population, definedness and causal restrictions. Exact targets and variants are editorial.
\subsection*{Definitions and precise research target}
\textit{Editorial formalization.} Consider low-income adults screening positive for depression in one district of India. Randomize one-time treatment offers \(z=(c,m)\in\{0,1\}^2\), where \(c\) offers a defined cash grant and \(m\) offers a defined course of evidence-based psychotherapy. For adult \(i\), let \(y_{it}(z)\) be real income and \(s_{it}(z)\) a validated symptom score at years \(t=1,3,5\). Define factorial interactions
\[\gamma_t^y=E[y_{it}(1,1)-y_{it}(1,0)-y_{it}(0,1)+y_{it}(0,0)],\]
and define \(\gamma_t^s\) by replacing income with symptoms. Estimate these interactions and all four treatment means under randomized offers, consistency, and measured village spillovers. A lasting positive income effect and lower symptoms may indicate reinforcing gains, but do not alone establish multiple stable equilibria. To evaluate a psychological poverty trap, specify a dynamic state-transition model and test whether its nonlinearities predict persistent divergence after temporary interventions, beyond baseline heterogeneity and regression to the mean. State attrition assumptions and collect consumption and work outcomes.

Define baseline depression eligibility with a fixed validated screening threshold before any assignment. Cash and therapy offers, eligibility, take-up, and therapy duration are fixed independently of post-treatment symptoms. Income is real annual income, with zeros permitted; symptom scores use one scale with larger values representing worse symptoms. Positivity and random assignment to all four arms identify factorial mean interactions under no interference. A poverty-trap model requires a state \(x_t=(y_t,s_t)\) and a declared transition kernel \(K(dx'\mid x,z_t)\); after temporary offers expire set \(z_t=0\). Multiple traps mean distinct stable invariant distributions or attracting states under that common untreated transition, with positive-probability basins. Finite-horizon mean effects alone do not prove such multiplicity. Attrition, mortality, and village exposure require explicit accounting. All expectations use a stated baseline population and probability law. Identification means every admissible data-generating model with the same observed-data law yields the same displayed target; otherwise report the identified set. Exact equations, horizons and assumptions are editorial, rather than source-given canonical conjectures.
\subsection*{Variants and assumption changes}
\begin{enumerate}
\item \textbf{Editorial sequencing variant.} Randomize cash-before-therapy versus therapy-before-cash at equal total doses and costs; define five-year income and symptom contrasts.
\item \textbf{Editorial state-threshold variant.} Specify \(x_{t+1}=f(x_t)+\xi_{t+1}\) after offers expire and randomize temporary dose near a candidate basin boundary. State finite-data bounds on persistent divergence and invariant states.
\end{enumerate}
\subsection*{References and source audit}
\textbf{1.} Explicit agenda on intervention combinations, long-run outcomes, and mental-health poverty traps. Locator: Science 370, eaay0214; full article pp.8–9 and Box 4. Full text or primary publication page inspected. URL: \url{https://economics.mit.edu/sites/default/files/inline-files/poverty-depression-anxiety-science.pdf}.
\begin{enumerate}\raggedright
\item\hypertarget{definitions:source:30007730:1}{} Matthew Ridley; Gautam Rao; Frank Schilbach; Vikram Patel. 2020. \emph{Poverty, depression, and anxiety: Causal evidence and mechanisms}. Science 370, eaay0214; full article pp.8–9 and Box 4.
\url{https://economics.mit.edu/sites/default/files/inline-files/poverty-depression-anxiety-science.pdf}.
DOI: \url{https://doi.org/10.1126/science.aay0214}.
\end{enumerate}

\subsection*{Common conventions: Source mathematical conventions}

These conventions apply to each entry unless explicitly replaced.

\textbf{Models and identification.} A primitive model \(m\) specifies all probability laws, feasible choices, information, equilibrium and measurement rules used by its target. The admissible class \(\mathcal M\) is fixed independently of outcomes. If \(P_O(m)\) is its observable probability law and \(\tau(m)\) the target, its identified set at observable law \(P\) is
\[
 \mathcal I_\tau(P)=\{\tau(m):m\in\mathcal M,\ P_O(m)=P\}.
\]
Point identification means that this set is a singleton. An empty set signals incompatibility of data and assumptions. Coordinate bounds need not describe the joint set. Bounds are sharp only if every retained target is attainable or arbitrarily approachable by compatible models. Finite bounds require explicit restrictions on unobserved quantities; unrestricted confounding, hidden wealth, extrapolation or arbitrary equilibrium selection can yield uninformative sets.

\textbf{Causal interventions.} A potential outcome \(Y_i(p)\) is the outcome under a complete policy regime \(p\), including timing, financing, information and market spillovers. The target population and units are fixed before treatment. With interference, use the complete assignment vector or a justified exposure mapping; individual no-interference notation alone cannot identify an economy-wide intervention. A difference of mean potential outcomes does not require the cross-world joint distribution. Conditioning on post-treatment migration, survival, enrollment or employment changes the estimand and needs separate justification.

\textbf{Statistical statements.} An estimator, confidence set, forecast or test specifies a sampling law, model class and loss/coverage criterion. Uniform \((1-\alpha)\)-coverage means \(\inf_{m\in\mathcal M}\Pr_m\{\tau(m)\in C_n\}\geq1-\alpha\), or an explicitly asymptotic version. The whole parameter space trivially has coverage, so informativeness requires a stated width, risk or power objective. A forecasting improvement is relative to a named benchmark, information set and evaluation population.

\textbf{Equilibrium and optimization.} An equilibrium comprises feasible plans, optimality, beliefs where needed, budget constraints and all market/resource clearing. The primitives or a stated benchmark define each correspondence \(\mathcal E(p;m)\). Nonemptiness is assumed or proved, never inferred just from compact policy sets. A selected-equilibrium target uses a declared selection rule; a robust target quantifies over all permitted equilibria. Use suprema or infima unless attainment is established. An \(\varepsilon\)-optimal policy is feasible and within \(\varepsilon\) of the finite optimum value. Report an empty feasible set separately.

\textbf{Welfare and accounting.} Normative utility, interpersonal normalization, social weights, numeraire, discounting and the use of public receipts are primitives. Behavioral utility need not coincide with the welfare criterion. Household welfare includes ownership income and rents once; adding profits again to compensating variation generally double-counts them. A monetary equivalent is defined by an explicit transfer and price convention, with monotonicity and range conditions. Average effects, mechanism contributions, transfers and welfare losses are different targets. Infinite sums require convergence and dynamic feasible plans require terminal or no-Ponzi conditions.

\textbf{Source versions.} A working-paper date, revision date and journal publication year are distinguished. A DOI for a journal version is not automatically the DOI of an earlier manuscript. Locators refer to the stated version and printed pages where specified. A metadata-only check does not verify a claim about a full-text conclusion. Every entry's source subsection narrows the provenance claim accordingly.
% END ECONOMICS STATEMENT
\end{document}
